% GENERATED FILE. Edit canonical lessons, then run npm run generate:books.
% canonical-source-hash: fnv1a64:f61d689577f8faee
% canonical-lessons: TE-C135-miriyalu, TE-C135-jilakarra, TE-C135-cintapandu, TE-C135-nimmakaya, TE-C135-draksa

\chapter{Black pepper, Cumin, Tamarind, Lemon, Grapes}
\label{ch:te-black-pepper-cumin-tamarind-lemon-grapes}
\hlchaptermodality{135}

\begin{chapteropening}
\textbf{By the end of this chapter:} \emph{I can say black pepper, cumin, tamarind, a lemon and grapes.}

Complete the last lesson of chapter 135: i can say black pepper, cumin, tamarind, a lemon and grapes.
\end{chapteropening}

\section[\texorpdfstring{\te{మిరియాలు} (miriyālu)}{మిరియాలు (miriyālu)}]{\te{మిరియాలు} (miriyālu) — black pepper}

\label{lesson:TE-C135-miriyalu}



\begin{quote}
Before the new one: say the Telugu for grain, then the Telugu for wheat.
\end{quote}

\subsection*{You'll want to know: \te{మిరియాలు}}
\textbf{\te{మిరియాలు}} — \emph{miriyālu} — \textquotedblleft{}black pepper\textquotedblright{}.

\begin{grammarlens}[title={What you've built}]
One more word for this chapter.
\end{grammarlens}

\subsection*{Your turn}
\begin{itemize}
\raggedright
  \item \emph{Say it:} \emph{miriyālu}
  \item \emph{Say it:} \emph{miriyālu}, once more
  \item \emph{Say it:} say \emph{miriyālu}
  \item \emph{Recall:} say the Telugu for grain, then the Telugu for wheat, then say \emph{miriyālu} again
\end{itemize}

\begin{tcolorbox}[breakable,colback=teal!4,colframe=teal!35!black,title={Before you move on}]
What does \te{మిరియాలు} mean? (\textquotedblleft{}Black pepper\textquotedblright{}.) Say it once more.
\end{tcolorbox}
\section[\texorpdfstring{\te{జీలకర్ర} (jīlakarra)}{జీలకర్ర (jīlakarra)}]{\te{జీలకర్ర} (jīlakarra) — cumin}

\label{lesson:TE-C135-jilakarra}



\begin{quote}
Before the new one: say the Telugu for wheat, then the Telugu for black pepper.
\end{quote}

\subsection*{You'll want to know: \te{జీలకర్ర}}
\textbf{\te{జీలకర్ర}} — \emph{jīlakarra} — \textquotedblleft{}cumin\textquotedblright{}.

\begin{grammarlens}[title={What you've built}]
One more word for this chapter.
\end{grammarlens}

\subsection*{Your turn}
\begin{itemize}
\raggedright
  \item \emph{Say it:} \emph{jīlakarra}
  \item \emph{Say it:} \emph{jīlakarra}, once more
  \item \emph{Say it:} say \emph{jīlakarra}
  \item \emph{Recall:} say the Telugu for wheat, then the Telugu for black pepper, then say \emph{jīlakarra} again
\end{itemize}

\begin{tcolorbox}[breakable,colback=teal!4,colframe=teal!35!black,title={Before you move on}]
What does \te{జీలకర్ర} mean? (\textquotedblleft{}Cumin\textquotedblright{}.) Say it once more.
\end{tcolorbox}
\section[\texorpdfstring{\te{చింతపండు} (cintapaṇḍu)}{చింతపండు (cintapaṇḍu)}]{\te{చింతపండు} (cintapaṇḍu) — tamarind}

\label{lesson:TE-C135-cintapandu}



\begin{quote}
Before the new one: say the Telugu for black pepper, then the Telugu for cumin.
\end{quote}

\subsection*{You'll want to know: \te{చింతపండు}}
\textbf{\te{చింతపండు}} — \emph{cintapaṇḍu} — \textquotedblleft{}tamarind\textquotedblright{}.

\begin{grammarlens}[title={What you've built}]
One more word for this chapter.
\end{grammarlens}

\subsection*{Your turn}
\begin{itemize}
\raggedright
  \item \emph{Say it:} \emph{cintapaṇḍu}
  \item \emph{Say it:} \emph{cintapaṇḍu}, once more
  \item \emph{Say it:} say \emph{cintapaṇḍu}
  \item \emph{Recall:} say the Telugu for black pepper, then the Telugu for cumin, then say \emph{cintapaṇḍu} again
\end{itemize}

\begin{tcolorbox}[breakable,colback=teal!4,colframe=teal!35!black,title={Before you move on}]
What does \te{చింతపండు} mean? (\textquotedblleft{}Tamarind\textquotedblright{}.) Say it once more.
\end{tcolorbox}
\section[\texorpdfstring{\te{నిమ్మకాయ} (nimmakāya)}{నిమ్మకాయ (nimmakāya)}]{\te{నిమ్మకాయ} (nimmakāya) — a lemon}

\label{lesson:TE-C135-nimmakaya}



\begin{quote}
Before the new one: say the Telugu for cumin, then the Telugu for tamarind.
\end{quote}

\subsection*{You'll want to know: \te{నిమ్మకాయ}}
\textbf{\te{నిమ్మకాయ}} — \emph{nimmakāya} — \textquotedblleft{}a lemon\textquotedblright{}.

\begin{grammarlens}[title={What you've built}]
One more word for this chapter.
\end{grammarlens}

\subsection*{Your turn}
\begin{itemize}
\raggedright
  \item \emph{Say it:} \emph{nimmakāya}
  \item \emph{Say it:} \emph{nimmakāya}, once more
  \item \emph{Say it:} say \emph{nimmakāya}
  \item \emph{Recall:} say the Telugu for cumin, then the Telugu for tamarind, then say \emph{nimmakāya} again
\end{itemize}

\begin{tcolorbox}[breakable,colback=teal!4,colframe=teal!35!black,title={Before you move on}]
What does \te{నిమ్మకాయ} mean? (\textquotedblleft{}A lemon\textquotedblright{}.) Say it once more.
\end{tcolorbox}
\section[\texorpdfstring{\te{ద్రాక్ష} (drākṣa)}{ద్రాక్ష (drākṣa)}]{\te{ద్రాక్ష} (drākṣa) — grapes}

\label{lesson:TE-C135-draksa}



\begin{quote}
Before the new one: say the Telugu for tamarind, then the Telugu for a lemon.
\end{quote}

\subsection*{You'll want to know: \te{ద్రాక్ష}}
\textbf{\te{ద్రాక్ష}} — \emph{drākṣa} — \textquotedblleft{}grapes\textquotedblright{}.

\begin{grammarlens}[title={What you've built}]
One more word for this chapter.
\end{grammarlens}

\subsection*{Your turn}
\begin{itemize}
\raggedright
  \item \emph{Say it:} \emph{drākṣa}
  \item \emph{Say it:} \emph{drākṣa}, once more
  \item \emph{Say it:} say \emph{drākṣa}
  \item \emph{Recall:} say the Telugu for tamarind, then the Telugu for a lemon, then say \emph{drākṣa} again
\end{itemize}

\begin{tcolorbox}[breakable,colback=teal!4,colframe=teal!35!black,title={Before you move on}]
What does \te{ద్రాక్ష} mean? (\textquotedblleft{}Grapes\textquotedblright{}.) Say it once more.
\end{tcolorbox}
